% Section 3. Sources: 602 PROOF.md sections 4-8 and REPORT.md sections 3-5; 1367 PROOF_CHECK.md (Lemma T),
% logs/check_r12.log.
%%%%%%%%%%%%%%%%%%%%%%%%%%%%%%%%%%%%%%%%%%%%%%%%%%%%%%%%%%%%%%%%%%%%%%%%
\section{The upper bound}\label{sec:upper}
%%%%%%%%%%%%%%%%%%%%%%%%%%%%%%%%%%%%%%%%%%%%%%%%%%%%%%%%%%%%%%%%%%%%%%%%

\subsection{A weight that the moves do not increase}

\begin{lemma}\label{lem:subinv}
Let $v>0$ be a function on all states and $R>0$, with
\[ \sum_{y\text{ successor of }x} v(y)\;\le\;R\,v(x)\qquad\text{for every state }x. \]
Then $z_n\le R^{2n}$ for every $n$.
\end{lemma}

\begin{proof}
By induction, the total weight of the endpoints of all walks of $t$ moves from $x$ is at most $R^tv(x)$. The closed
walks at $e$ contribute $(A^{2n})_{e,e}\,v(e)$ to this total for $t=2n$; divide by $v(e)$.
\end{proof}

The bound holds for every $n$ with constant one, with no asymptotic input. The work is to find $v$ and to check
infinitely many inequalities with finitely many computations.

\subsection{The weight}

Fix a \emph{radius} $r\ge2$. For a stack $S$ let $E(S)$ be the number of adjacent equal pairs (positions $i$ with
$S_i=S_{i+1}$) and $T(S)$ the number of adjacent equal triples. Let $\mathrm{view}(S)$ be the top $r$ symbols of $S$,
or all of $S$ together with a bottom mark if $S$ is shorter. The weight is
\[ v(U,L,c)\;=\;a^{|U|+|L|}\,p^{E(U)+E(L)}\,t^{T(U)+T(L)}\;u\bigl(\mathrm{view}(U),\mathrm{view}(L),c\bigr), \]
with three positive numbers $a,p,t$ and a table $u$ of positive integers indexed by the pairs of views and the
parity. Intuitively, $a$ prices the number of open chords, $p$ and $t$ price runs of equal parities, which are rare
in typical stacks, and the table corrects the price for the top of each stack, where the moves act.

The inequality of Lemma~\ref{lem:subinv} at a state depends only on its top $r+1$ symbols in each stack.

\begin{lemma}[Truncation]\label{lem:trunc}
Let $x'$ be obtained from $x$ by keeping the top $r+1$ symbols of each stack. Then
$\sum_y v(y)/v(x)=\sum_{y'}v(y')/v(x')$.
\end{lemma}

\begin{proof}
The allowed moves agree, because they depend only on the top symbols. For each move the quotient $v(y)/v(x)$ is
$a^{\pm1}$ times $p$ and $t$ to the changes of $E$ and $T$, times $u(y)/u(x)$. A push or pop changes $E$ and $T$ only
through the top three symbols, and the views before and after the move are read from the top $r+1$ symbols; these
agree for $x$ and $x'$ since $r+1\ge3$.
\end{proof}

So the infinitely many inequalities reduce to the states with at most $r+1$ symbols in each stack:
$2(2^{r+2}-1)^2$ of them, which is $536{,}805{,}378$ at $r=12$. In the language of~\cite{PaperB} these are the rows of
the certificate; here each row is simply a short state.

\subsection{The certificate}

The certificate has radius $12$, parameters
\[ a=\tfrac{1783}{2000},\qquad p=\tfrac{137}{200},\qquad t=\tfrac{97}{125},\qquad R=\tfrac{3187011595}{10^9}, \]
and a table of $67{,}092{,}481$ used entries (the entries for parity $1$ are those for parity $0$ with the two
parities exchanged, which is only a way of storing the table). All inequalities hold; with $R$ lowered by
$10^{-9}$ four of them fail, so $R$ is the least admissible value with this denominator. The largest quotient
$\sum_yv(y)/v(x)$ is $3.187011594348$, at $U=0$, $L=1111000000$ (bottom first), $c=0$. By
Lemma~\ref{lem:subinv}, $z_n\le R^{2n}$, which is the upper half of Theorem~\ref{thm:main}.

The table was found by an iterative solver in floating point; the solver plays no role in the proof. Smaller radii
give weaker bounds: $R^2=10.187205938372711824$ at radius $7$, and in floating point $10.2646$, $10.2323$,
$10.2135$, $10.1978$ at radii $3$ to $6$. Without the factors $p$ and $t$ the bound at radius $8$ is $10.2963$, worse
than radius $3$ with them. The gap to the estimate $10.113$ shrinks roughly like $r^{-1}$; a gap of $0.01$ would
need a radius of about $50$, far beyond tables of size $4^{r+1}$.
